% !TEX program = xelatex
\documentclass[a4paper]{article}
\usepackage{amsthm}
\usepackage{amssymb}
\usepackage{bm}
\usepackage{mathtools}
\usepackage[x11names]{xcolor}
\usepackage{xparse}
\usepackage{fontspec}
\usepackage{unicode-math}
\setromanfont[Ligatures={Common}]{DovesType-Regular.otf}
\setsansfont{Andika}
\setmathfont{Asana Math}[Scale=1]
\newfontfamily\TrajanP[Scale=1.0]{TRAJANPRO-REGULAR.OTF}
\newfontfamily{\cnfamily}[Scale=1.0]{SIMSUN.TTC}
% \usepackage{CJKutf8}
% \usepackage{pstricks}
\usepackage{varwidth}
\usepackage{siunitx}
\usepackage{graphicx}
\usepackage[margin=1.5cm,top=2.0cm]{geometry}
\usepackage[most]{tcolorbox}
\usepackage{pgfplots}
\pgfplotsset{compat=newest}
\tcbuselibrary{skins,xparse,poster,breakable}
% \usetikzlibrary{fadings}
\usetikzlibrary{calc,math, plotmarks, shapes, shapes.geometric, positioning, angles, intersections, quotes, through, patterns, turtle, arrows.meta}
\usetikzlibrary{decorations.markings,backgrounds}
% \usepackage{etoolbox}
\usepackage{tkz-euclide}
% \usepackage{xlop}
% \newcommand\hole[2]{#1}  % for use with xlop
\usepackage{enumitem}
\pagenumbering{gobble}
%%%%%%%%%%%%%%%%%%%%%%%%%%%%%%%%%%%%%%%%%%%%%%%%%%%%%%%%%
\newcommand\markangle[9]{% origin X Y radius radiusmark mark colour opacity
%  % fill red circle offset-from-centre
  \begin{scope}
    \path[clip] (#1) -- (#2) -- (#3);
    \fill[color=#7,fill opacity=#8,draw=black,name path global=pcircle]  % global declaration required otherwise pcircle is not seen by the `named intersections=' lines below.
    (#1) circle (#4);
  \end{scope}
  % middle calculation
  \path[name path=line one] (#1) -- (#2);
  \path[name path=line two] (#1) -- (#3);
  \path[%
  name intersections={of=line one and pcircle, by={inter one}},
  name intersections={of=line two and pcircle, by={inter two}}
  ] (inter one) -- (inter two) coordinate[pos=#9] (place);
  % put mark
  \node at ($(#1)!#5!(place)$) {\scriptsize{#6}};
}

\newcommand\tcircle[6]{% centre coord (x,y), radius, points, radpoint, colour, edge
  \coordinate (O) at (#1,#2); % centre of the circle
  \def\radius{#3}          % radius of the circle
  \def\npts{#4}            % number of the points
  \def\radpt{#5}           % radius of the points
  \colorlet{ptcolour}{#6}  % colour of the points
  % \draw (O) circle (\radius);
  \foreach \numpoint in {1,...,\npts}{
    \fill[ptcolour] (O) ++ (360/\npts*\numpoint:\radius) coordinate (C\numpoint) circle(\radpt);
  }
}

% \newcommand{\condSoln}[2]{\ifcsdef{r@#1}{#2}{}}

% \newcommand\fadingtext[3][]{%
%    \begin{tikzfadingfrompicture}[name=fading letter]
%      \node[text=transparent!0,inner xsep=0pt,outer xsep=0pt,#1] {#3};
%    \end{tikzfadingfrompicture}%
%    \begin{tikzpicture}[baseline=(textnode.base)]
%      \node[inner sep=0pt,outer sep=0pt,#1](textnode){\phantom{#3}};
%      \shade[path fading=fading letter,#2,fit fading=false]
%      (textnode.south west) rectangle (textnode.north east);%
%    \end{tikzpicture}%
% }

\definecolor{JISpurple}{RGB}{89,72,122}
\definecolor{JISivory}{RGB}{241,234,221}
\definecolor{JIStaupe}{RGB}{183,156,154}
\definecolor{PaleGreen}{RGB}{240,255,240} % 'Honeydew'

\AddToHook{shipout/background}{%
    \put (0in,-\paperheight){\includegraphics[width=\paperwidth,height=\paperheight]{images/cnybkg1.jpg}}%
}

\newcommand\numberthis{\addtocounter{equation}{1}\tag{\theequation}}

\newtcolorbox{MyOuterBox}{%
  enhanced,
  watermark graphics=images/CNYrabbitWM.jpg,
  watermark opacity=0.8,
  watermark overzoom=1.0,
  breakable,
  frame style=JISpurple,
  colback=JISivory,
  colframe=JISpurple,
  title={\includegraphics[width=0.9cm,height=0.9cm]{images/JIS Final Logo FA-02.png}\raisebox{3mm}{\Large{CNY Maths Challenge}\hspace{22em} \Large{\bfseries\sffamily 18}}},
}

\newtcolorbox{MyInnerBox}[2][]{enhanced,%empty,
coltitle=JISpurple,colback=white,
breakable,
%left skip=2.5cm,
fonttitle=\bfseries\sffamily,
attach boxed title to top left={yshift=-1.5mm},
boxed title style={empty, size=small, top=1mm, bottom=0pt},
varwidth boxed title=0.5\linewidth,
frame code={
  \path (title.east|-frame.north) coordinate (aux);
\path[draw=JISpurple, line width=0.5mm, rounded corners,fill=white]
(frame.west) |- ([xshift=-2.5mm]title.north east) to[out=0, in=180] ([xshift=7.5mm]aux)-|(frame.east)|-(frame.south)-|cycle;
},
title={#2},#1}

\newtcolorbox{MyInnerSplitBox}[2][]{enhanced,%empty,
bicolor,sidebyside,sidebyside align=top seam,
righthand width=7.5cm,colbacklower=white,
sidebyside gap=5mm,
breakable,
coltitle=JISpurple,colback=white,
fonttitle=\bfseries\sffamily,
attach boxed title to top left={yshift=-1.5mm},
boxed title style={empty, size=small, top=1mm, bottom=0pt},
varwidth boxed title=0.5\linewidth,
frame code={
  \path (title.east|-frame.north) coordinate (aux);
\path[draw=JISpurple, line width=0.5mm, rounded corners,fill=white]
(frame.west) |- ([xshift=-2.5mm]title.north east) to[out=0, in=180] ([xshift=7.5mm]aux)-|(frame.east)|-(frame.south)-|cycle;
},
title={#2},#1}


\newtcolorbox{MySolutionBox}[1][]{%
  enhanced,
  breakable,
  frame style=JISpurple,
  colback=PaleGreen, colframe=green,
  title={\Large Solution},
  drop fuzzy shadow,
  halign=left,
  #1
}

%%%%%%%%%%%%%%%%%%%%%%%%%%%%%%%%%%%%%%%%%%%%%%%%%%
\newtoggle{SOLUTION}
%%% Uncomment the appropriate line below to show solutions %%%
\toggletrue{SOLUTION}
% \togglefalse{SOLUTION}
%%%%%%%%%%%%%%%%%%%%%%%%%%%%%%%%%%%%%%%%%%%%%%%%%


%%%%%%%%%%%%%%%%%%%%%%%%%%%%%%%%%%%%%%%%%%%%%%%%%%
%%%%%%            DOCUMENT BEGINS           %%%%%%
%%%%%%%%%%%%%%%%%%%%%%%%%%%%%%%%%%%%%%%%%%%%%%%%%%
\begin{document}
  \begin{MyOuterBox}
    \iftoggle{SOLUTION}{Here are the full, or partial solutions.
    }{
      Welcome to this week's Maths Challenge!\\
      Have a go at both questions!\\
      Drop your solution in the box in the staffroom by Tuesday
    }
       \begin{MyInnerBox}{Year 9 and below}
      A game begins with 100 people standing in a circle. One person holds a ball, call them person Number 1, and then number the other players consecutively anticlockwise: the person to the ball-holder's left is No. 100, to their right stands person No. 2.\par
      The person holding the ball throws it over the head of the first person on their right, the person on the other side catches it. So at the start No. 1 throws to No. 3.\par
      A person skipped by the ball, immediately leaves the circle for tea and chocolate biscuits.\par
      So person No. 3 would be the next holder of the ball. No. 3 throws to No. 5 and No. 4 leaves the circle.\par
      The process continues until there is one person left in the circle. If there are only two people left, the person holding the ball throws to themselves, eliminating the other person.\par
      What is the number of the last person in the circle?\par
      \iftoggle{SOLUTION}{%conditional output begin
      \begin{MySolutionBox}
      \underline{Solution by playing the full game:} Each line of the table shows the people left in the circle after that round:\par
        \begin{description}[leftmargin=!,labelwidth=\widthof{\bfseries Round 7:},labelindent=1em,labelsep=2em,itemsep=0.1pt]
          \item[\textbf{Round 1:}] 1, 3, 5, 7, 9, 11, 13, 15, 17, 19, 21, 23, 25, 27, 29, 31, 33, 35, 37, 39, 41, 43, 45, 47, 49, 51, 53, 55, 57, 59, 61, 63, 65, 67, 69, 71, 73, 75, 77, 79, 81, 83, 85, 87, 89, 91, 93, 95, 97, 99
          \item[Round 2:] 1, 5, 9, 13, 17, 21, 25, 29, 33, 37, 41, 45, 49, 53, 57, 61, 65, 69, 73, 77, 81, 85, 89, 93, 97
          \item[Round 3:] 1, 9, 17, 25, 33, 41, 49, 57, 65, 73, 81, 89, 97
          \item[Round 4:] 9, 25, 41, 57, 73, 89
          \item[Round 5:] 9, 41, 73
          \item[Round 6:] 9, 73
          \item[Round 7:] 73
        \end{description}
        So person No. 73 is the last one in the circle.\par
        \underline{Solution by reasoning:} For an even number of people starting the game, half are out after the first round.\par
        So if the number of people playing is a power of \(2\), the person holding the ball at the beginning will be the last person in the circle.\par
        The closest power of \(2\), less than \(100\) is \(64\). So if we can see who is holding the ball once \(36\) people have been eliminated, we have found the answer.\par
        The \(36^{t}\) person to be eliminated must be No. 72, and No. 73 has just caught the ball.\par
        Since there are now \(64\) people in the circle, the holder of the ball, No. 73, will be the last person in the circle.\par
      \end{MySolutionBox}
    }{}%conditional output end
    \end{MyInnerBox}


    \vspace{0.0cm}
        \begin{MyInnerBox}{Year 10 and above}
       One hundred people take part in a race in which no-one can tie. Later, after the race, every competitor is asked about what position they came, and all answer, giving a number between 1 and 100. The sum of all their answers is 4000. What is the minimum number of people who have lied about what position in the race they achieved?\par
      \iftoggle{SOLUTION}{%conditional output begin
      \begin{MySolutionBox}
        If all the competitors told the truth about their final position in the race then the sum of their answers should be \(1+2+3+\cdots +98+99+100\). This is an arithmetic sequence with first term \(a=1\) and constant difference \(d=1\). We can use the formula for the sum of the first \(n\) digits.
        \begin{align*}
          S_{n}&=\frac{n}{2}\left(2a+(n-1)d\right)\\
          S_{100}&=\frac{100}{2}\left(2\times 1 + (100-1)\times 1\right)\\
          S_{100}&=50\times(2+99) = 5050
      \end{align*}
        Or we could notice, as Gauss did (aged 7!), that rearranging our sum:
        \begin{align*}
          (1+100) + (2+99) +\cdots+ (49+52) + (50+51) &= 101 + 101 + \cdots + 101 + 101\\
          \shortintertext{(\(50\) pairs of digits, each pair adding to \(101\))}
                                                      &= 101\times 50\\
        &=5050
        \end{align*}
        The scores obtained from the runners is only \(4000\) though, there is \(5050-4000=1050\) worth of fibbing going on!\par
        Now, for the minimum number of people lying, we want the 'missing' \(1050\) to be accounted for by the smallest number of people, so we choose the people who came \(100^{th},\; 99^{th}\) and so on.\par
        Each of these people has to report a made up position, to minimise the number of liars, the number they give should be as small as possible, so we'll assume they all say they came first!\par
        Let's see what's happening:\par\medskip
        \begin{tabular}{l|l|l|l|l|l}
          Person & Real & False & Amount & Total & Total\\
          number & place & place & reduced & before & after\\\hline
          \(1\) & \(100\) & \(1\) & \(99\) & \(5050\) & \(4951\)\\
          \(2\) & \(99\) & \(1\) & \(98\) & \(4951\) & \(4853\)\\
          \(3\) & \(98\) & \(1\) & \(97\) & \(4853\) & \(4756\)\\
          \(\vdots\) & \(\vdots\) & \(\vdots\) & \(\vdots\) & \(\vdots\)\\
          \(11\) & \(90\) & \(1\) & \(89\) & \(4105\) & \(4016\)\\ \hline
           & & & \(1034\) & & \\
        \end{tabular}\par\medskip
        Eleven competitors lying has reduced the total to \(4016\), we need one more person to lie about their finishing place to make the total of scores \(4000\). If the person in \(17^{th}\) place says they came first, then the sum of scores is \(4000\) as required.\par
        So twelve people is the minimum number who lied.\par
      \end{MySolutionBox}
    }{}%conditional output end
    \end{MyInnerBox}


  \end{MyOuterBox}

%%%%%%%%%%%%%%%%%%%%%%%%%%%%%%%%%%%%%%%%%%%%%%%%%%
\end{document}



